\begin{abstract}
% Reconstructing a deformable target moving over a large area is challenging. Existing methods for the 3D reconstruction of a moving articulated object require dense coverage in both the viewing space and the temporal dimension. Therefore they typically rely on videos captured from multiple viewpoints at each time step. 
Reconstructing a dynamic target moving over a large area is challenging. Standard approaches for dynamic object reconstruction require dense coverage in both the viewing space and the temporal dimension, typically relying on multi-view videos captured at each time step.
However, such setups are only possible in constrained environments.
In real-world scenarios, observations are often sparse over time and captured sparsely from diverse viewpoints (e.g., from security cameras), making dynamic reconstruction highly ill-posed. 
We present \ourmethod, a framework that simultaneously estimates a deformation model and the object’s motion over time under sparse observations.
To initialize \ourmethod, we
leverage a rough skeleton graph and an initial static reconstruction as inputs to guide motion estimation. (Later, we show that this input requirement can be relaxed.)
Our method optimizes a skeleton-driven deformation field composed of a coarse skeleton joint pose estimator and a module for fine-grained deformations.
% , per-bone influence radii, a skinning correction field, and a fine detail deformation field.
By making only the joint pose estimator time-dependent, our model enables smooth motion interpolation while preserving learned geometric details. 
Experiments on synthetic datasets show that our method outperforms existing approaches under sparse observations by up to 34\% in PSNR, and achieves comparable performance to dense monocular video methods on real-world datasets despite using significantly fewer frames. Moreover, we demonstrate that the input initial static reconstruction can be replaced by a diffusion-based generative prior, making our method more practical for real-world scenarios. 


\end{abstract}